
\documentclass{article}
\usepackage{colortbl}
\usepackage{makecell}
\usepackage{multirow}
\usepackage{supertabular}

\begin{document}

\newcounter{utterance}

\centering \large Interaction Transcript for game `wordle', experiment `medium\_frequency\_words\_no\_clue\_no\_critic', episode 4 with reflexion{-}Qwen3.5{-}9B.
\vspace{24pt}

{ \footnotesize  \setcounter{utterance}{1}
\setlength{\tabcolsep}{0pt}
\begin{supertabular}{c@{$\;$}|p{.15\linewidth}@{}p{.15\linewidth}p{.15\linewidth}p{.15\linewidth}p{.15\linewidth}p{.15\linewidth}}
    \# & \multicolumn{2}{c}{Player} && \multicolumn{2}{c}{Game Master} \\
    \hline

    \theutterance \stepcounter{utterance}  
    & & & \multicolumn{4}{p{0.6\linewidth}}{
        \cellcolor[rgb]{0.9,0.9,0.9}{
            \makecell[{{p{\linewidth}}}]{
                \texttt{\tiny{[P1$\langle$GM]}}
                \texttt{You have attempted this same task before and did not succeed. These are the plans you wrote after those attempts.} \\
\texttt{Trial \#1: New plan: The last attempt was aborted because "star" is only four letters long, violating the rule that the target word must be five lowercase letters. Although my explanation correctly identified that 'r', 'a', and 'e' are present and misplaced, I failed to construct a valid five{-}letter word by simply appending 't' to "star". Next time, I will ensure my new guess is exactly five letters long by combining the confirmed letters ('r', 'a', 'e') with two new, untested letters (excluding 'c' and 'n') to form a valid English word, such as "arise" or "raise", rather than creating a truncated string.} \\
\\ 
\\ 
\texttt{You are a language wizard who likes to guess words by using the given rules.} \\
\\ 
\texttt{Welcome to Wordle! You have six attempts to guess the target word, a valid English word of five lowercase letters (a{-}z). Please use the tags "explanation:" and "guess:" to provide a concise explanation for each guess.} \\
\\ 
\texttt{For instance, if your guess is "apple", your response should be} \\
\texttt{explanation: this is a common five{-}letter English word, and I am starting my guess with this word.} \\
\texttt{guess: apple} \\
\\ 
\texttt{After each guess, your answer will be validated, and you will receive feedback indicating which letters are correct (green), which letters are correct but in the wrong position (yellow), and which letters are incorrect (red). This feedback can be useful in determining which letters to include or exclude in your next guess.} \\
\\ 
\texttt{For example, the feedback for "apple" might be:} \\
\texttt{guess\_feedback: a<yellow> p<yellow> p<green> l<yellow> e<red>} \\
\\ 
\texttt{The explanation should contain details about how the guess\_feedback is used to arrive at a new guess.} \\
\\ 
\texttt{Do not generate any other text.} \\
\\ 
\texttt{Let's begin with your first guess.} \\
            }
        }
    }
     \\ \\

    \theutterance \stepcounter{utterance}  
    & \multicolumn{4}{p{0.6\linewidth}}{
        \cellcolor[rgb]{0.561,0.737,0.561}{
            \makecell[{{p{\linewidth}}}]{
                \texttt{\tiny{[P1$\rangle$GM]}}
                \texttt{explanation: I am starting with a high{-}frequency word containing common vowels and consonants to maximize information gain in the first attempt.} \\
\texttt{guess: crane} \\
            }
        }
    }
    & & \\ \\

    \theutterance \stepcounter{utterance}  
    & & & \multicolumn{2}{p{0.3\linewidth}}{
        \cellcolor[rgb]{0.9,0.9,0.9}{
            \makecell[{{p{\linewidth}}}]{
                \texttt{\tiny{[GM$|$GM]}}
                \texttt{attempts = 1} \\
\texttt{target\_word = lager} \\
\texttt{guess = crane} \\
\texttt{guess\_feedback = c<red> r<yellow> a<yellow> n<red> e<yellow>} \\
            }
        }
    }
    & & \\ \\

    \theutterance \stepcounter{utterance}  
    & & & \multicolumn{4}{p{0.6\linewidth}}{
        \cellcolor[rgb]{0.9,0.9,0.9}{
            \makecell[{{p{\linewidth}}}]{
                \texttt{\tiny{[P1$\langle$GM]}}
                \texttt{guess\_feedback: c<red> r<yellow> a<yellow> n<red> e<yellow>} \\
\\ 
\texttt{Provide your response only in this format.} \\
\texttt{explanation: details} \\
\texttt{guess: word} \\
            }
        }
    }
     \\ \\

    \theutterance \stepcounter{utterance}  
    & \multicolumn{4}{p{0.6\linewidth}}{
        \cellcolor[rgb]{0.561,0.737,0.561}{
            \makecell[{{p{\linewidth}}}]{
                \texttt{\tiny{[P1$\rangle$GM]}}
                \texttt{explanation: The feedback indicates that 'r', 'a', and 'e' are present but misplaced, while 'c' and 'n' are incorrect. Since 'a' is yellow in the second position and 'e' is yellow in the last position, I will try placing 'a' in the first position and 'e' in the third position, keeping 'r' in the second position to form "arete" (a valid word meaning a peak or summit), though I suspect "arise" or "raise" might be better fits if 'r' moves. However, given the constraints, I will test "arise" to see if 'r' belongs at the start or if 'i' is the missing vowel. Wait, 'r' is yellow in the second spot in "crane", so in "arise", 'r' is second. 'a' is first (was second in crane), 'i' is new, 's' is new, 'e' is last (was last in crane). This tests new letters 'i' and 's' while repositioning knowns.} \\
\texttt{guess: arise} \\
            }
        }
    }
    & & \\ \\

    \theutterance \stepcounter{utterance}  
    & & & \multicolumn{2}{p{0.3\linewidth}}{
        \cellcolor[rgb]{0.9,0.9,0.9}{
            \makecell[{{p{\linewidth}}}]{
                \texttt{\tiny{[GM$|$GM]}}
                \texttt{attempts = 2} \\
\texttt{target\_word = lager} \\
\texttt{guess = arise} \\
\texttt{guess\_feedback = a<yellow> r<yellow> i<red> s<red> e<yellow>} \\
            }
        }
    }
    & & \\ \\

    \theutterance \stepcounter{utterance}  
    & & & \multicolumn{4}{p{0.6\linewidth}}{
        \cellcolor[rgb]{0.9,0.9,0.9}{
            \makecell[{{p{\linewidth}}}]{
                \texttt{\tiny{[P1$\langle$GM]}}
                \texttt{guess\_feedback: a<yellow> r<yellow> i<red> s<red> e<yellow>} \\
\\ 
\texttt{Provide your response only in this format.} \\
\texttt{explanation: details} \\
\texttt{guess: word} \\
            }
        }
    }
     \\ \\

    \theutterance \stepcounter{utterance}  
    & \multicolumn{4}{p{0.6\linewidth}}{
        \cellcolor[rgb]{0.561,0.737,0.561}{
            \makecell[{{p{\linewidth}}}]{
                \texttt{\tiny{[P1$\rangle$GM]}}
                \texttt{explanation: The feedback shows 'a', 'r', and 'e' are present but misplaced, while 'i' and 's' are incorrect. Since 'a' was yellow in the first position of "arise" and 'r' was yellow in the second, and 'e' is yellow in the last, I need to rearrange them. 'a' cannot be first, 'r' cannot be second, and 'e' cannot be last. Trying "rarely" is too long. Let's try "reave" (to tear away) or "reale" (not a word). How about "reave"? No, 'v' is new. Let's try "reale" is invalid. What about "reave"? r{-}e{-}a{-}v{-}e. 'r' first (was 2nd), 'e' second (was 3rd), 'a' third (was 1st), 'v' new, 'e' last (was last). But 'e' is already last and yellow, so maybe it stays? Wait, if 'e' is last and yellow, it could be correct or wrong position. In "arise", 'e' was last and yellow. If I keep 'e' last, I need to place 'a' and 'r' elsewhere. 'a' cannot be first. 'r' cannot be second. Try "reave": r{-}e{-}a{-}v{-}e. 'r' first, 'e' second, 'a' third, 'v' fourth, 'e' fifth. But 'e' is repeated. Is "reave" valid? Yes. But maybe "reale" isn't. How about "reave"? Or "reare"? "reare" is not common. Let's try "reave". Wait, 'v' is untested. Maybe "reale" is not a word. What about "reave"? r{-}e{-}a{-}v{-}e. Feedback: r (was 2nd, now 1st), e (was 3rd, now 2nd), a (was 1st, now 3rd), v (new), e (was 5th, now 5th). If 'e' is correct at the end, it should be green. But in "arise", 'e' was yellow at the end, meaning it's present but not at the end? Or it was yellow because it was in the wrong spot? In "arise", 'e' is at the end. If it's yellow, it means 'e' is in the word but not at the end. So 'e' cannot be last. Ah! That's the key. In "arise", 'e' was at the end and got yellow, so 'e' is NOT at the end. So I must move 'e' away from the last position. '} \\
            }
        }
    }
    & & \\ \\

    \theutterance \stepcounter{utterance}  
    & & & \multicolumn{2}{p{0.3\linewidth}}{
        \cellcolor[rgb]{0.9,0.9,0.9}{
            \makecell[{{p{\linewidth}}}]{
                \texttt{\tiny{[GM$|$GM]}}
                \texttt{Error: The response should contain the 'guess:' keyword exactly once.} \\
            }
        }
    }
    & & \\ \\

    \theutterance \stepcounter{utterance}  
    & & & \multicolumn{2}{p{0.3\linewidth}}{
        \cellcolor[rgb]{0.9,0.9,0.9}{
            \makecell[{{p{\linewidth}}}]{
                \texttt{\tiny{[GM$|$GM]}}
                \texttt{game\_result = ABORT} \\
            }
        }
    }
    & & \\ \\

    \theutterance \stepcounter{utterance}  
    & & \multicolumn{4}{p{0.6\linewidth}}{
        \cellcolor[rgb]{0.95,0.95,0.95}{
            \makecell[{{p{\linewidth}}}]{
                \texttt{\tiny{[P1$\leftrightarrow$P]}}
                \texttt{Post{-}episode Reflexion (not sent to GM):} \\
\texttt{New plan: The last attempt was aborted because I generated a multi{-}sentence explanation containing internal monologue and contradictory reasoning (e.g., debating "reave" vs. "reare") instead of a single, concise explanation followed immediately by the guess. The rules strictly require the response to contain only the tags "explanation:" and "guess:" with no other text. Additionally, I failed to correctly deduce that 'e' cannot be in the last position since it was yellow in "arise". Next time, I will strictly adhere to the format: provide a brief, single{-}sentence explanation of the deduction (specifically noting 'e' is not last and 'a'/'r' are misplaced), and immediately follow with a valid five{-}letter word that repositions these letters, such as "rare" (invalid length) {-}> "rarely" (invalid length) {-}> "rears" or "rares" or "reave" (checking constraints). A better candidate based on 'a', 'r', 'e' present but misplaced, and 'i', 's' absent: "reave" (r{-}e{-}a{-}v{-}e) puts 'e' in the middle, but 'e' was yellow at the end in "arise", so 'e' is not at the end. "rears" (r{-}e{-}a{-}r{-}s) has 'e' not at end. Wait, 'r' is repeated. Let's try "reave" again? No, 'e' at end is bad. How about "reare"? No. "reave" is r{-}e{-}a{-}v{-}e. 'e' is at 2 and 5. If 'e' is not at 5, "reave" fails. Try "rears": r{-}e{-}a{-}r{-}s. 'e' is 2nd. 'a' is 3rd. 'r' is 1st and 4th. This works if 'r' can be 1st. But 'r' was yellow at 2nd in "crane". In "arise", 'r' was yellow at 2nd. So 'r' is not 2nd. In "rears", 'r' is 1st and 4th. This is valid. But I need to ensure the word is valid. "Rears" is valid. Or "reave" is invalid due to 'e' position. Let's pick "rears" or "rares". Actually, "reave" is r{-}e{-}a{-}v{-}e. If 'e' is not last, "reave" is bad. "rears" is good. I will construct a response with exactly one explanation sentence and one guess line, ensuring the guess is a valid 5{-}letter word that respects the feedback (a, r, e present but not in current spots; i, s absent; c, n absent;} \\
            }
        }
    }
    & \\ \\

\end{supertabular}
}

\end{document}
